\documentclass[11pt,a4paper]{article}

% ---------------------------------------------------------
% PREAMBLE (stable, warning-free, Overleaf-safe)
% ---------------------------------------------------------
\usepackage[utf8]{inputenc}
\usepackage[T1]{fontenc}
\usepackage{lmodern}
\usepackage{amsmath, amssymb, amsfonts}
\usepackage{bm}
\usepackage{geometry}
\usepackage{graphicx}
\usepackage{hyperref}
\usepackage{cite}
\usepackage{authblk}
\usepackage{physics}
\usepackage{mathtools}

\geometry{margin=2.4cm}

\title{\textbf{Companion LXX: Wavelet Transforms of Persistence Diagrams in the Relaxation-Driven Cyclic Cosmology}}
\author{Michael Lehmann \\ \texttt{mi.lehmann@gmx.de}}
\date{April 2026}

\begin{document}
\maketitle

\begin{abstract}
Wavelet transforms of persistence diagrams provide a joint scale–position 
representation of topological fluctuations in cosmological fields. In weak-lensing 
analyses, wavelet coefficients quantify how the birth and death of connected 
components, loops, and void-like structures vary across both scale and filtration 
parameter. In the standard cosmological model, the wavelet structure reflects the 
nonlinear matter distribution and the projection of large-scale structure. In the 
Relaxation-Driven Cyclic Cosmology (RDCC), the scale-dependent evolution of the 
gravitational potential modifies the multi-scale topology in a characteristic way. 
This Companion develops a continuous description of wavelet transforms in RDCC, 
deriving how the relaxation scale influences multi-scale topological coherence, 
wavelet-energy distributions, and the observational signatures in weak-lensing 
surveys. The resulting framework provides a distinctive probe of the RDCC 
gravitational field through multi-scale topology.
\end{abstract}
% ---------------------------------------------------------
\section{Introduction}

Persistence diagrams encode the birth and death of topological features in 
weak-lensing convergence fields. While spectral density functions provide a 
frequency-resolved view of topological fluctuations, wavelet transforms offer a 
joint representation in scale and position, capturing localized multi-scale 
structure.

In the standard cosmological model, the wavelet structure reflects the nonlinear 
matter distribution and the projection of large-scale structure. In RDCC, the 
gravitational potential evolves differently. The relaxation mechanism suppresses 
the decay of small-scale modes and enhances the coherence of large-scale modes. 
This modifies the multi-scale topology in a scale-dependent manner, producing a 
distinctive signature in wavelet analyses.
% ---------------------------------------------------------
\section{Wavelet Transforms of Persistence Diagrams}

Given a filtration-ordered sequence of persistence diagrams $\{D(\nu)\}$, the 
continuous wavelet transform (CWT) of the Wasserstein-distance fluctuations is

\begin{equation}
    W_{D}(a, b)
    = \int
      W_{p}(D(\nu), D^{*})
      \psi^{*}\!\left(\frac{\nu - b}{a}\right)
      \frac{d\nu}{\sqrt{|a|}},
\end{equation}

where $a$ is the scale, $b$ is the position, and $\psi$ is the mother wavelet.

The wavelet energy is

\begin{equation}
    E_{D}(a)
    = \int |W_{D}(a, b)|^{2} \, db.
\end{equation}
% ---------------------------------------------------------
\section{RDCC Modification of Persistence Diagrams}

In RDCC, the convergence evolves as
\begin{equation}
    \kappa_{\mathrm{RDCC}}(\ell)
    = \kappa(\ell)
      \left[
        1 - \chi_{\kappa}
        \left(\frac{\ell}{\ell_{\mathrm{rel}}}\right)^{\alpha}
      \right].
\end{equation}

This modifies the birth and death thresholds of topological features:

\begin{align}
    b_{\mathrm{RDCC}} &= b \left[ 1 - \chi_{b} (\ell/\ell_{\mathrm{rel}})^{\alpha} \right], \\
    d_{\mathrm{RDCC}} &= d \left[ 1 - \chi_{d} (\ell/\ell_{\mathrm{rel}})^{\alpha} \right].
\end{align}

Thus the RDCC wavelet coefficients become

\begin{equation}
    W_{D,\mathrm{RDCC}}(a, b)
    = W_{D}(a, b)
      \left[
        1 - \chi_{W}
        \left(\frac{\ell}{\ell_{\mathrm{rel}}}\right)^{\alpha}
      \right].
\end{equation}
% ---------------------------------------------------------
\section{Wavelet Energy in RDCC}

The RDCC wavelet energy is

\begin{equation}
    E_{D,\mathrm{RDCC}}(a)
    = E_{D}(a)
      \left[
        1 - \chi_{E}
        \left(\frac{\ell}{\ell_{\mathrm{rel}}}\right)^{\alpha}
      \right].
\end{equation}

This modification reflects the relaxation-driven change in multi-scale topology.
% ---------------------------------------------------------
\section{Multi-Scale Topological Signatures of RDCC}

The relaxation mechanism enhances the coherence of large-scale modes. Wavelet 
transforms therefore exhibit:

\begin{itemize}
    \item enhanced wavelet energy at large scales,
    \item suppressed wavelet energy at small scales,
    \item a characteristic turnover near $\ell_{\mathrm{rel}}$,
    \item modified scale–position localization of topological features,
    \item altered alignment of persistence lifetimes across scales.
\end{itemize}

These signatures provide a direct probe of the RDCC gravitational field through 
multi-scale topology.
% ---------------------------------------------------------
\section{Conclusion}

Wavelet transforms of persistence diagrams in the Relaxation-Driven Cyclic 
Cosmology reflect the scale-dependent evolution of the gravitational potential and 
the nonlinear matter distribution. The relaxation mechanism modifies the multi-
scale topology in a scale-dependent manner, producing distinctive signatures in 
weak-lensing surveys. These effects provide a direct observational window into the 
relaxation scale and the temporal structure of the RDCC gravitational field. The 
present Companion establishes the theoretical foundation for future studies of 
multi-scale topology, nonlinear structure, and the interplay between light and 
gravity in RDCC.

\bibliographystyle{apsrev4-2}
\nocite{*}
\bibliography{references}


\end{document}
